\documentclass[11pt,a4paper]{article}
% ---------------------------------------------------------------------------
%  FermiWick -- common style of the showcase and of the manual.
%  The examples are computed by Development/run_examples.wl: examples/NAME.wl is the code, NAME.tex the output as a
%  formula, NAME.txt the output as text.
% ---------------------------------------------------------------------------
\usepackage[utf8]{inputenc}
\usepackage[T1]{fontenc}
\usepackage[sc]{mathpazo}
\usepackage[scaled=0.88]{helvet}
\renewcommand{\ttdefault}{txtt}
\usepackage{amsmath,amssymb,bm}
\allowdisplaybreaks[2]
\usepackage[a4paper,margin=2.2cm,top=2.4cm,bottom=2.4cm,headheight=22pt]{geometry}
\usepackage{microtype}
\usepackage[dvipsnames]{xcolor}
\usepackage{graphicx}
\usepackage{listings}
\usepackage{simpler-wick}
\usepackage[most]{tcolorbox}
\usepackage{booktabs,array,tabularx}
\usepackage{titlesec}
\usepackage{enumitem}
\usepackage{fancyhdr}
\usepackage{needspace}
\usepackage{etoolbox}
\usepackage[colorlinks=true,linkcolor=accent,urlcolor=accent,citecolor=accent]{hyperref}

\definecolor{accent}{HTML}{8E2A7B}
\definecolor{accentlight}{HTML}{F1DCEC}
\definecolor{accentpale}{HTML}{FAF0F7}
\definecolor{codebg}{HTML}{FBF6FA}
\definecolor{codeframe}{HTML}{E1CCDC}
\definecolor{keyword}{HTML}{8E2A7B}
\definecolor{comment}{HTML}{7D6A78}
\definecolor{checkgreen}{HTML}{2E7D32}

\titleformat{\section}{\sffamily\Large\bfseries\color{accent}}{\thesection}{0.8em}{}
\titleformat{\subsection}{\sffamily\large\bfseries\color{accent}}{\thesubsection}{0.7em}{}
\titleformat{\paragraph}[runin]{\sffamily\bfseries\color{accent}}{}{0pt}{}[.]
\titlespacing*{\section}{0pt}{2.6ex plus 1ex minus 0.4ex}{1.2ex}
\titlespacing*{\subsection}{0pt}{2.2ex plus 0.6ex minus 0.3ex}{0.8ex}
\titlespacing*{\paragraph}{0pt}{1.2ex plus 0.4ex}{0.6em}
\makeatletter
\pretocmd{\section}{\needspace{8\baselineskip}}{}{}
\pretocmd{\subsection}{\if@nobreak\else\needspace{6\baselineskip}\fi}{}{}
\makeatother
\setcounter{tocdepth}{2}
\setlist{itemsep=1pt, topsep=3pt}
\setlength{\emergencystretch}{1.5em}
\setlength{\parindent}{0pt}
\setlength{\parskip}{4pt plus 1.5pt minus 1pt}

\pagestyle{fancy}
\fancyhf{}
\fancyhead[R]{\sffamily\small\color{comment}\nouppercase{\rightmark}}
\fancyfoot[C]{\sffamily\small\thepage}
\renewcommand{\sectionmark}[1]{\markboth{\thesection\enspace #1}{\thesection\enspace #1}}
\renewcommand{\headrulewidth}{0.4pt}
\renewcommand{\headrule}{\hbox to\headwidth{\color{codeframe}\leaders\hrule height \headrulewidth\hfill}}

\lstdefinelanguage{WL}{
  morekeywords={NormalOrder,OperatorCommutator,OperatorAnticommutator,ToSpinChargeBasis,ToLadderBasis,GroupLocalOperators,
    OperatorCoefficients,CleanDeltas,ResetNormalOrder,DefineOperator,QuasiFreeState,SpinSingletState,GeneralState,StateSelect,
    LocalPairQ,AnomalousPairQ,Contraction,WickVEV,VerifyNormalOrder,MeanField,DefinePeierlsHopping,DefineFlux,
    GutzwillerProjector,LocalSandwich,GutzwillerVEV,GutzwillerSimplify,SdotS,SdotScrossS,SpinOperator,EmptyProjector,
    SingleOccupationProjector,cDag,nOp,nTotOp,SzOp,SpOp,SmOp,DOp,Bdag,
    FockSpace,FockSpaceQ,FockSector,FockSectorQ,ModeIndex,OccupationBasis,BasisState,StateToKet,FockMatrix,FockExpectation,
    SpinSquared,ParticleHoleOperator,TimeReversal,SectorBlock,LowestStates,LehmannGreen,GreenMatrix,ClearFockCache},
  sensitive=true, morecomment=[n]{(*}{*)}, morestring=[b]",
}
\lstset{language=WL, basicstyle=\ttfamily\small, keywordstyle=\color{keyword}\bfseries, commentstyle=\color{comment}\itshape,
  stringstyle=\color{BrickRed}, columns=fullflexible, keepspaces=true, breaklines=true, breakatwhitespace=true,
  breakindent=2em, showstringspaces=false, upquote=true, extendedchars=true,
  literate={·}{{$\cdot$}}1 {φ}{{$\varphi$}}1 {Φ}{{$\Phi$}}1 {σ}{{$\sigma$}}1 {ρ}{{$\rho$}}1 {κ}{{$\kappa$}}1 {↑}{{$\uparrow$}}1 {↓}{{$\downarrow$}}1 {χ}{{$\chi$}}1 {Δ}{{$\Delta$}}1 {θ}{{$\theta$}}1 {ν}{{$\nu$}}1 {η}{{$\eta$}}1}
\lstdefinestyle{plainoutput}{language={}, basicstyle=\ttfamily\small, breaklines=true, columns=fullflexible, keepspaces=true}

\newcommand{\pkg}[1]{\texttt{#1}}
\newcommand{\fn}[1]{\texttt{\color{accent}#1}}
\newcommand{\ok}{\textcolor{checkgreen}{\ensuremath{\checkmark}}}
\newcommand{\exmissing}[1]{\textcolor{red}{[missing example: \detokenize{#1}]}}
\makeatletter
\newcommand{\excodeafter}{\tcb@parfillskip@check\par\ifvmode\nobreak\@nobreaktrue\nointerlineskip\addvspace{\glueexpr2pt-\parskip}\fi}
\makeatother
\newcommand{\excode}[1]{\ifhmode{\widowpenalty=10000\par}\nobreak\fi\IfFileExists{examples/#1.wl}{\tcbinputlisting{listing only,
  colback=codebg, colframe=codeframe, boxrule=0.5pt, arc=2pt, left=6pt, right=6pt, top=3pt, bottom=3pt, before skip=4pt,
  after=\excodeafter, listing file={examples/#1.wl}, listing options={language=WL}}}{\exmissing{#1}}}
\newtcolorbox{exoutput}{colback=white, colframe=codeframe, boxrule=0.5pt, arc=2pt, left=24pt, right=6pt, top=3pt, bottom=3pt,
  before skip=2pt, after skip=8pt, breakable, overlay unbroken and first={\node[anchor=north west,
  font=\scriptsize\sffamily\color{comment}] at ([xshift=4pt, yshift=-5pt]frame.north west) {out};}}
\newcommand{\exformula}[1]{\begin{exoutput}\raggedright\IfFileExists{examples/#1.tex}{$\displaystyle\input{examples/#1.tex}$}{\exmissing{#1}}\end{exoutput}}
\newcommand{\extext}[1]{\IfFileExists{examples/#1.txt}{\tcbinputlisting{listing only, colback=white, colframe=codeframe,
  boxrule=0.5pt, arc=2pt, left=24pt, right=6pt, top=2pt, bottom=2pt, before skip=2pt, after skip=8pt,
  listing file={examples/#1.txt}, listing options={style=plainoutput}, overlay unbroken and first={\node[anchor=north west,
  font=\scriptsize\sffamily\color{comment}] at ([xshift=4pt, yshift=-5pt]frame.north west) {out};}}}{\exmissing{#1}}}
\newcommand{\example}[1]{\excode{#1}\IfFileExists{examples/#1.tex}{\exformula{#1}}{\IfFileExists{examples/#1.txt}{\extext{#1}}{}}}

\newcommand{\titleblock}[2]{%
  \begin{center}
  \includegraphics[width=2.2cm]{../Logo/FermiWick_logo.pdf}\\[6pt]
  {\sffamily\Huge\bfseries\color{accent}FermiWick}\\[4pt]
  {\sffamily\Large #1}\\[8pt]
  {\small #2}\\[2pt]
  {\small Nazzareno Africani (SISSA) \textperiodcentered{} version 0.2.0 \textperiodcentered{}
   \href{https://github.com/n4sd4q/FermiWick}{github.com/n4sd4q/FermiWick}}
  \end{center}\vspace{6pt}}
\newtcblisting{code}{listing only, colback=codebg, colframe=codeframe, boxrule=0.5pt, arc=2pt, left=6pt, right=6pt,
  top=3pt, bottom=3pt, listing options={language=WL}}
\newcommand{\extypeset}[1]{\begin{exoutput}\centering\small\IfFileExists{examples/#1_typeset.tex}{$\displaystyle\input{examples/#1_typeset.tex}$}{\exmissing{#1}}\end{exoutput}}

\familysetup{document=Showcase}

\begin{document}
\familytitle{Operators, contractions, mean field, Gutzwiller projectors and matrices}{Everything the package does, step by
step: every output is computed by the code printed above it}{version 0.3.0 \textperiodcentered{} November 2026}

\begin{tcolorbox}[colback=accentpale, colframe=accentpale, arc=2pt, left=8pt, right=8pt]
\small FermiWick is a Wolfram Language package for the algebra of fermion ladder operators. This document goes through
all it does, in the order of a calculation: the algebra of the operators (normal order, local operators, the
spin--charge basis and the recognition of spin operators); the contractions of a Gaussian state, declared by rules or
taken from ready-made states, and Wick's theorem; the mean-field Hamiltonian of an operator of any order; Gutzwiller
projectors; and the same operators as matrices, for exact diagonalisation. The code of each section uses only what the
previous ones introduced.
\end{tcolorbox}

{\small\tableofcontents}

% =====================================================================================================================
\section{The algebra of the operators}

\subsection{Ladder operators and normal order}
Fermions are $c_{i\sigma}$ and $c^\dagger_{i\sigma}$, written \fn{c[i, s]} and \fn{cDag[i, s]}. The site $i$ is any expression;
the spin $\sigma$ is $\uparrow$, $\downarrow$ or a symbol. Products are written with \texttt{**}, in the order given, and
\fn{NormalOrder} moves every $c^\dagger$ to the left with $\{c_{i\sigma},c^\dagger_{j\sigma'}\}=\delta_{ij}\delta_{\sigma\sigma'}$:
\example{car}
Two different symbolic sites are different sites, so their operators anticommute; a symbolic spin is split into $\uparrow$
and $\downarrow$:
\example{distinctsites}
\example{samesite}
Two equal fermions annihilate the product, wherever they are:
\example{pauli}
The product $a\cdot b$ (\texttt{\textbackslash[CenterDot]}) is the normal-ordered product. \fn{OperatorCommutator} and
\fn{OperatorAnticommutator} are normal-ordered too:
\example{commutator1}
\example{anticommutator}

\subsection{Local operators, kept compact}
The local operators of a site are $n_{i\sigma}$, $n_i$, $S^z_i$, $S^\pm_i$, the double occupancy
$D_i=n_{i\uparrow}n_{i\downarrow}$, the doublon operators $B^\dagger_i=c^\dagger_{i\uparrow}c^\dagger_{i\downarrow}$ and $B_i$, and the
projectors $P^0_i$ (empty) and $P^1_i$ (singly occupied); their definitions:
% The composite operators of FermiWick and their definitions (included by the showcase and by the manual).
\begin{center}\small
\renewcommand{\arraystretch}{1.25}
\begin{tabular}{lll}\toprule
FermiWick & operator & definition \\ \midrule
\texttt{nOp[i, s]} & $n_{i\sigma}$ & $c^\dagger_{i\sigma}c_{i\sigma}$ \\
\texttt{nTotOp[i]}, \texttt{nOp[i]} & $n_i$ & $n_{i\uparrow}+n_{i\downarrow}$ \\
\texttt{SzOp[i]} & $S^z_i$ & $\tfrac12\,(n_{i\uparrow}-n_{i\downarrow})$ \\
\texttt{SpOp[i]}, \texttt{SmOp[i]} & $S^+_i$, $S^-_i$ & $c^\dagger_{i\uparrow}c_{i\downarrow}$, \ $c^\dagger_{i\downarrow}c_{i\uparrow}$ \\
\texttt{SpinOperator[i, a]} & $S^a_i$ & $\tfrac12\sum_{\sigma\sigma'}c^\dagger_{i\sigma}\,\sigma^a_{\sigma\sigma'}\,c_{i\sigma'}$, \ $a=1,2,3$ \\
\texttt{DOp[i]} & $D_i$ (double occupancy) & $n_{i\uparrow}n_{i\downarrow}$ \\
\texttt{Bdag[i]}, \texttt{B[i]} & $B^\dagger_i$, $B_i$ (doublon) & $c^\dagger_{i\uparrow}c^\dagger_{i\downarrow}$, \ $c_{i\downarrow}c_{i\uparrow}$ \\
\texttt{EmptyProjector[i]} & $P^0_i$ & $(1-n_{i\uparrow})(1-n_{i\downarrow})=1-n_i+D_i$ \\
\texttt{SingleOccupationProjector[i]} & $P^1_i$ & $n_i-2D_i$ \\
\texttt{GutzwillerProjector[i]} & $P^G_i$ & $\sqrt2\,[\cos\theta\,P^1_i+\sin\theta\,(P^0_i+D_i)]$ \\
\texttt{SdotS[i, j]} & $\mathbf S_i\cdot\mathbf S_j$ & $S^z_iS^z_j+\tfrac12(S^+_iS^-_j+S^-_iS^+_j)$ \\
\texttt{SdotScrossS[i, j, k]} & $\mathbf S_i\cdot(\mathbf S_j\times\mathbf S_k)$ & $\sum_{abc}\epsilon_{abc}\,S^a_iS^b_jS^c_k$ \\ \bottomrule
\end{tabular}
\end{center}
Products of these on one site reduce with the algebra of the fermions (e.g.\ $D_in_i=2D_i$, $(S^z_i)^2=\tfrac14n_i-\tfrac12D_i$);
operators on several sites and products with symbolic sites stay as products. \texttt{GutzwillerProjector} is in
\texttt{FermiWick`Gutzwiller`}; one-site operators on several sites are products, \texttt{DOp[i, j]} $=D_iD_j$.

They are written compactly and expanded only when needed; \fn{ToLadderBasis} expands them:
\example{ladder}
\fn{ToSpinChargeBasis} goes the other way: it writes a product with the local operators of each site. A product is
separated site by site, the sites being reordered with the fermion sign only, and each site is reduced once:
\example{compact}
On one site the algebra of the local operators follows from that of the fermions: $S^+S^-=n_\uparrow-D$,
$(S^z)^2=\frac14n-\frac12D$, and $\mathbf S^2=\frac34(n-2D)$, a spin $\frac12$ only on singly occupied sites:
\example{spinflip}
\example{szsz}
\example{spinlength}
The projectors are complete, and $P^1$ commutes with the spin:
\example{projectors}
\example{projectorproduct}
\fn{GroupLocalOperators} keeps the ladder operators but groups them site by site, with the fermion sign:
\example{group}

\subsection{The spin--charge basis and the recognition of spin operators}
A product of ladder operators on several sites is decomposed into local operators; spin flips become $S^+_1S^-_2$:
\example{decompose}
\fn{ToSpinChargeBasis} then recognizes the operators that act on several sites. The sum of the products of the spin
components is the exchange $\mathbf S_1\cdot\mathbf S_2$, and \texttt{"Recognize" -> False} shows its expansion:
\example{recognizess}
\example{norecognize}
The recognition allows a remainder in the diagonal part, so that an XXZ coupling is read as a Heisenberg term plus an
Ising correction:
\example{xxz}
The commutator of two exchange terms that share a site is the scalar chirality,
$[\mathbf S_1\cdot\mathbf S_2,\mathbf S_2\cdot\mathbf S_3]=-i\,\mathbf S_1\cdot(\mathbf S_2\times\mathbf S_3)$; the chirality is
antisymmetric in its sites and is written with them in order; and the Heisenberg triangle commutes with the total spin:
\example{commutator}
\example{chiralityorder}
\example{invariance}
These are identities of the whole Fock space, not only of the singly occupied sector.

\subsection{Operators of your own}
\fn{DefineOperator} defines an operator on several sites by its expression in known operators; it stays compact, is
expanded when needed and recognized in the results. A spin-conserving hopping on a bond, and its commutator with the
density at its end:
\example{define}
The recognition keeps the fermion signs, also with the sites out of order: the Hermitian conjugate of the hopping from
1 to 2 is the hopping from 2 to 1.
\example{definedagger}

\subsection{The Hermitian conjugate}
\fn{Dagger} reverses the products, exchanges $c$ and $c^\dagger$ (and $S^+$, $S^-$; $B$, $B^\dagger$) and conjugates the
coefficients. The amplitudes $\chi$, $\Delta$, $\rho$, $\kappa$ are complex; the other symbols, couplings, angles and phases,
are real unless declared complex with \texttt{"Complex" -> \{\dots\}}:
\example{dagger}
\fn{PlusHC} adds the Hermitian conjugate, the ``$+$ h.c.'' of a Hamiltonian, and \fn{HermitianOperatorQ} checks it:
\example{plushc}
\example{hermitianq}

\subsection{Modes with any labels}
A mode can carry any labels: \fn{c[l1, l2, \dots]}. \fn{DeclareModeLabels} says what each label is: \texttt{"Site"} (used for
locality; different symbolic sites are different sites), \texttt{"Spin"} (a spin $\frac12$, which enables the spin operators),
or any other name, a generic label compared with Kronecker deltas. By default \fn{c[i]} is a spinless fermion on site $i$,
\fn{c[i, s]} has a site and a spin, and more labels are a site and generic labels. Spinless fermions:
\example{spinless}
Orbitals: the anticommutator gives a delta of the orbitals, kept symbolic:
\example{orbitals}
Momenta, labels with no site: the deltas stay, and multiply what they act on:
\example{momenta}

\subsection{Coefficients and lattice phases}
\fn{OperatorCoefficients} gives the coefficient of each monomial in normal order:
\example{coefficients}
\fn{ToTeX} writes any expression in \LaTeX:
\example{totex}
\fn{DefinePeierlsHopping} and \fn{DefineFlux} attach phases to the hoppings and fix the flux through a loop:
\example{peierls}

% =====================================================================================================================
\section{Contractions and Wick's theorem}
A Gaussian state is fixed by its contractions $\langle O_1O_2\rangle$, and Wick's theorem reduces every expectation value
to them: \fn{WickVEV} normal-orders the expression and sums over the pairings of each monomial, with the sign of the
permutation. Such states, Gaussian in the fermions, are also called quasi-free. In FermiWick the state is an object, declared in one of two ways: by its symmetries, which say which
contractions vanish, or by rules for the contractions.

\subsection{The general state and its symmetries}
On a pair of sites the normal contractions form a $2\times2$ matrix in spin, and so do the anomalous ones. Written in
spin channels,
\[
\langle c^\dagger_{i\sigma}c_{j\sigma'}\rangle=\chi_{ij}\,\delta_{\sigma\sigma'}+\textstyle\sum_a\chi^a_{ij}\,\sigma^a_{\sigma'\sigma},\qquad
\langle c_{i\sigma}c_{j\sigma'}\rangle=\big[(\Delta_{ij}+\textstyle\sum_a\Delta^a_{ij}\,\sigma^a)\,\epsilon\big]_{\sigma\sigma'},\qquad
\epsilon_{\uparrow\downarrow}=1,
\]
with the charge channels $\chi$ (hopping per spin) and $\Delta$ (singlet pairing, symmetric in the sites) and the spin
channels $\chi^a$ and $\Delta^a$ (the triplet $d$ vector, antisymmetric in the sites), $a=x,y,z$. Hermiticity,
$\chi_{ji}=\chi_{ij}^*$ and $\chi^a_{ji}=\chi^{a*}_{ij}$, always holds. On one site $\chi_{ii}$ is the density per spin,
$\chi^a_{ii}=\langle S^a_i\rangle$ the magnetization and $\Delta_{ii}=\langle c_{i\uparrow}c_{i\downarrow}\rangle$. \fn{GaussianState[]}
is this general state; \fn{ContractionTable} shows its contractions:
\example{generallocal}
\example{generalpairing}
A symmetry sets channels to zero or makes them real, and the state is declared by its symmetries:
\begin{center}\small
\begin{tabular}{ll}\toprule
symmetry & what it leaves \\ \midrule
\texttt{"SpinRotation"} (SU(2)) & the charge channels $\chi$, $\Delta$ \\
\texttt{"SpinRotationZ"} (rotations about $z$) & $\chi$, $\chi^z$, $\Delta$, $\Delta^z$ \\
\texttt{"Charge"} (U(1) of the particle number) & no pairing \\
\texttt{"TimeReversal"} & $\chi$, $\Delta$ real; $\chi^a$, $\Delta^a$ imaginary; no magnetization \\ \bottomrule
\end{tabular}
\end{center}
With $S^z$ conserved the pairing is a singlet plus the $m=0$ triplet; with SU(2) and time reversal the hopping is real
and symmetric:
\example{symz}
\example{symsu2tr}
The exchange in a state that conserves $S^z$ and the particle number has the classical term of the magnetizations, the
hopping and the $z$ spin channel:
\example{szexchange}

\subsection{States declared by rules}
\fn{GaussianState} takes rules for the contractions, with patterns in the operators. Pairs without a rule contract to
zero, and $\langle c\,c^\dagger\rangle$ follows from $\langle c^\dagger c\rangle$ by the anticommutator. Two sites with
hopping $g$ and half filling:
\example{rulesstate}
\example{contraction}
The rules can have any spin structure: here a hopping $h$ that flips the spin, besides the hopping $g$ that keeps it;
it enters the correlation of $S^z$ with the opposite sign:
\example{flipstate}

\subsection{States invariant under spin rotations}
\fn{SpinSingletState} is the most common state, written short: invariant under spin rotations, with hopping
$\chi_{ij}=\langle c^\dagger_{i\sigma}c_{j\sigma}\rangle$, filling $\frac12$ per spin and, on request, singlet pairing
$\Delta_{ij}=\langle c_{i\uparrow}c_{j\downarrow}\rangle$. Without pairing it is \texttt{GaussianState[\{"SpinRotation", "Charge"\}, "Filling" -> 1/2]}:
\example{singletequal}
\example{singletstate}
In the exchange, hopping and pairing enter with the same weight, as the SU(2) gauge structure of spins written with
fermions requires~\cite{AZHA88}:
\example{exchange}
The chirality is a closed loop of three contractions, three hoppings or one hopping and two pairings; its imaginary part is
the flux of the amplitudes through the triangle:
\example{chirality}
Filling and on-site pairing are options:
\example{doubleocc}

\subsection{Magnetic and general states}
A state need not be invariant. Local magnetizations $m_i$, $\langle n_{i\sigma}\rangle=\frac12\pm m_i$, with no
contractions between sites, give the classical exchange:
\example{magnetic}
\example{magneticexchange}
The same state follows from the symmetry: rotations about $z$, the magnetization $m_i$ on each site, and only the local
contractions:
\example{magneticsym}
\fn{GeneralState} has normal contractions $\rho_{i\sigma j\sigma'}$ of any spin structure, here a transverse magnetization,
and anomalous ones $\kappa$ on request:
\example{general}
\example{generalanomalous}

\subsection{Selecting and truncating contractions}
Symmetries describe states; filters select terms. \fn{StateSelect} keeps the contractions that satisfy a criterion: the
result is a physical state when the criterion is a symmetry or a decoupling of the sites, and otherwise a way to isolate a
class of terms. Without the anomalous ones the pairing state is a
pure-hopping state; with only the local ones the sites are uncorrelated; with only the non-local ones the Hartree term is
gone:
\example{select}
\example{localonly}
\example{nonlocalonly}
\texttt{"MaxOrder" -> m} keeps the terms with at most $m$ contractions between different sites, and
\fn{VerifyNormalOrder} checks Wick's theorem with and without normal ordering:
\example{maxorder}
\example{verify}

\subsection{Wick's theorem written out}
\fn{WickPairings} lists the pairings of a product that do not vanish, with their signs and values; \fn{WickTeX} writes the
whole expansion in \LaTeX, with the contractions drawn by the package \pkg{simpler-wick} (two \LaTeX{} runs). The pairings
of a hopping back and forth, at half filling:
\example{pairings}
Hopping and pairing in a product of four operators: the string, and the string typeset:
\example{wicktex}
\extypeset{wicktex}
Only the non-local contractions of a loop around a triangle; overlapping contractions are drawn at different heights:
\example{wicktexloop}
\extypeset{wicktexloop}

% =====================================================================================================================
\section{Mean field}
\subsection{The mean-field Hamiltonian of any operator}
The Hartree--Fock--Bogoliubov Hamiltonian is the quadratic operator whose expectation value has the same first variation
as $\mathcal E=\langle H\rangle$ among Gaussian states. For a product of $2n$ fermions it is obtained by taking two
operators out in every possible way, in their order and with the sign of the permutation that brings them to the front,
and replacing the rest by its expectation value; the constant $-(n-1)\langle O_1\cdots O_{2n}\rangle$ gives
$\langle H_{\rm MF}\rangle=\mathcal E$. \fn{MeanField} does this for any order and any state. The Hubbard interaction gives
Hartree and BCS terms in the singlet state, and a spin-dependent Hartree field in the magnetic one:
\example{hubbard}
\example{magnetichubbard}
A density--density interaction has Hartree and Fock terms; a quadratic operator is its own mean field:
\example{density}
\example{quadratic}

\subsection{The fields of a Hamiltonian}
\fn{OperatorCoefficients} reads the fields. For the exchange $J\,\mathbf S_1\cdot\mathbf S_2$: the hopping field
$t_{12}=-\frac34J\chi^*_{12}$, the pairing field $D_{12}=\frac34J\Delta_{12}$ and the constant:
\example{exchangefields}
The chirality has six fermions, so its fields are quadratic in the amplitudes: each bond of the triangle receives a
hopping from the path through the third site, and a pairing from one hopping and one pairing on the other two sides:
\example{chiralfields}

\subsection{Hermiticity and the variational property}
The mean-field Hamiltonian of a Hermitian operator is Hermitian, and has the expectation value of the operator,
identically:
\example{hermitianmf}
\example{variational}

% =====================================================================================================================
\section{Gutzwiller projectors}\label{sec:gutzwiller}
\pkg{FermiWick`Gutzwiller`} adds the half-filling Gutzwiller operator
$P^G_i=\sqrt2\,[\cos\theta\,P^1_i+\sin\theta\,(P^0_i+D_i)]$, with $\theta=0$ the projection on single occupation, and
\fn{GutzwillerVEV}, the expectation value to first order in the contractions between different sites. In the
spin--charge basis:
\example{gutzwillerprojector}
At half filling $P^G$ is normalized for every $\theta$, and the double occupancy is $d=\frac12\sin^2\theta$:
\example{gutzwillernorm}
\example{gutzwillerd}
In a state with real hopping (time reversal), the hopping between two projected sites is reduced by
\example{hopping}
which is the factor $q=8d(1-2d)$ of Brinkman and Rice~\cite{BR70}:
\example{brq}
Minimizing $8d(1-2d)\,\varepsilon_0+Ud$, with $\varepsilon_0<0$ the kinetic energy of the free electrons per bond, gives
\example{brd}
so that $d$ vanishes at $U_c=8|\varepsilon_0|$, the Brinkman--Rice transition. The exchange, instead, is enhanced: for
$\theta=0$ it is four times its value in the free state, the factor $g_s=4$ of Zhang, Gros, Rice and Shiba at half
filling~\cite{ZGRS88}:
\example{projectedexchange}

% =====================================================================================================================
\section{Matrices and exact diagonalisation}
\pkg{FermiWick`Fock`} writes the operators of the previous sections as sparse matrices on the Fock space of a few sites,
so that a cluster can be solved exactly and a symbolic result checked. It is loaded, in a cell of its own, by
\example{fockload}
A space is given by its sites (and orbitals, or no
spin); its modes are the ladder operators, numbered site by site with the spin fastest, and a basis state is the string of
their occupations, the first mode being the most significant bit:
\example{fockspace}
A sector keeps the states with a given particle number and $S^z$, and every operator is built directly in it, without
the whole space:
\example{focksector}
\fn{FockMatrix} reads any expression of the package: ladder, local and defined operators, products in the order written,
symbolic coefficients. The Hubbard dimer, written as on paper, has in this sector the spectrum
\example{dimerspectrum}
The eigenvectors are exact as well, and \fn{StateToKet} writes a state with the occupations of the sites. The ground
state, the lowest of the four energies for all $t,U>0$ (the numbers only rank them), is a singlet whose weight on the
doubly occupied sites falls as $U$ grows:
\example{dimerground}
At a finite temperature, the free energy of the whole Fock space, sixteen states, at the chemical potential $U/2$ of
half filling:
\example{dimerfree}
The entropy $S=-\partial F/\partial T$ goes from zero, the ground state being a singlet, to $\log 16$, the sixteen states
equally likely (at $U=4t$ for the first limit):
\example{dimerentropy}
The Green's function at zero temperature, analytically: its Lehmann sum is written with the matrices of
$c^\dagger_{1\uparrow}$ and $c_{1\uparrow}$ from the sector of the ground state to those with three and one particles,
\fn{FockMatrix}\texttt{[op, \{from, to\}]}, and with the eigenstates of those sectors:
\example{dimergreen}
Its four poles, at $\pm\big(\frac12\sqrt{16t^2+U^2}\pm t\big)$, are the energies of adding and of removing an electron.

Ground states come from \fn{LowestStates} (dense, or Arnoldi on large sectors) and expectation values from
\fn{FockExpectation}. On a ring of eight sites at half filling with $U=4t$ (a sector of dimension 4900), the energy per
site, the double occupancy and the correlation of neighbouring spins are
\example{ring8}
\fn{LehmannGreen} gives the Green's function at zero temperature in the Lehmann representation; the poles of the dimer at
$U=4t$, with the ground state at half filling, are at the energies of adding or removing a particle:
\example{lehmann}
Wick's theorem is exact in a Gaussian state, and here it can be checked: the contractions of the free ground state of the
ring with six electrons (a closed shell, so that the ground state is a single Slater determinant), computed from the
matrices, define a \fn{GaussianState}, and \fn{WickVEV} of an interacting operator in that state is its exact
expectation value:
\example{wickexact}

% =====================================================================================================================
\section{Checks}
The test suite (\texttt{Tests/RunTests.wls}) checks the algebra, the mean-field formulas of the exchange, of the chirality
and of the Hubbard interaction, the states, and the identities of this showcase. It checks the symbolic engine against the
matrices: the matrix of an expression equals that of its normal order and of its spin--charge form, and \fn{WickVEV} and
\fn{MeanField} give the exact values in a random Gibbs state of a quadratic Hamiltonian with pairing, on three sites with
spin and on two sites with two orbitals. It checks the matrices themselves against the Jordan--Wigner construction, and
exact diagonalisation on the Hubbard dimer, a magnetic plaquette and Hubbard rings: the residuals are below $10^{-10}$.

\begin{thebibliography}{9}\small
\bibitem{AZHA88} I.~Affleck, Z.~Zou, T.~Hsu, and P.~W.~Anderson, Phys.\ Rev.\ B {\bf 38}, 745 (1988).
\bibitem{BR70} W.~F.~Brinkman and T.~M.~Rice, Phys.\ Rev.\ B {\bf 2}, 4302 (1970).
\bibitem{ZGRS88} F.~C.~Zhang, C.~Gros, T.~M.~Rice, and H.~Shiba, Supercond.\ Sci.\ Technol.\ {\bf 1}, 36 (1988).
\end{thebibliography}
\end{document}
